\documentclass[../../../include/open-logic-section]{subfiles}

\begin{document}

\olfileid{tur}{mac}{una}
\olsection{数的一元表示}

% Part: turing-machines% Chapter: machines-computations% Section: unary-numbers

\begin{explain}
图灵机处理写在其纸带上的符号序列。
根据图灵机所使用的字母表不同，这些符号序列可以表示各种输入和输出。
当然，我们特别关注的是那些计算\emph{算术}函数（即自然数上的函数）的图灵机。
表示正整数的一种简单方法，是将它们编码为由单个符号~$\TMstroke$ 组成的序列。
若 $n \in \Nat$，当 $n = 0$ 时，令 $\TMstroke^n$ 为空序列；否则，令其为恰好由 $n$ 个 $\TMstroke$ 组成的序列。
\end{explain}

\begin{defn}[计算]
图灵机~$M$ \emph{计算}函数 $f\colon \Nat^k \to \Nat$，当且仅当
$M$ 在输入
\[
\TMstroke^{n_1} \TMblank \TMstroke^{n_2} \TMblank \dots \TMblank \TMstroke^{n_k}
\]
上停机且输出 $\TMstroke^{f(n_1, \dots, n_k)}$。
\end{defn}

\begin{prob}
  请给出图灵机~$M$ 计算函数 $f\colon \Nat^k \to \Nat^m$ 的定义。
\end{prob}

\begin{ex}\ollabel{ex:adder}
\emph{加法：}
我们来构造一台计算函数 $f(n,m) = n + m$ 的机器。
这要求机器启动时纸带上有两个长度分别为 $n$ 和~$m$ 的 $\TMstroke$ 块，
而停机时纸带上有一个由 $n+m$ 个~$\TMstroke$ 组成的块。
这两个输入的 $\TMstroke$ 块由一个 $\TMblank$ 分隔，
因此一种方法是在包含 $\TMblank$ 的方格上写一个 $\TMstroke$，并擦除最后一个 $\TMstroke$。
\begin{figure}\[
\begin{tikzpicture}[->,>=stealth',shorten >=1pt,auto,node distance=2.8cm,
                    semithick]
  \tikzstyle{every state}=[fill=none,draw=black,text=black]

  \node[initial,state]         (A)              {$q_0$};
  \node[state]         (B) [right of=A] {$q_1$};
  \node[state]         (C) [right of=B] {$q_2$};

  \path (A) edge node {\TMtrans{\TMblank}{\TMstroke}{\TMstay}} (B)
           edge [loop above] node {\TMtrans{\TMstroke}{\TMstroke}{\TMright}} (B)
        (B) edge [loop above] node {\TMtrans{\TMstroke}{\TMstroke}{\TMright}} (B)
            edge node {\TMtrans{\TMblank}{\TMblank}{\TMleft}} (C)
        (C) edge [loop above] node {\TMtrans{\TMstroke}{\TMblank}{\TMstay}} (C);
\end{tikzpicture}
\]\caption{计算 $f(x,y) = x+y$ 的机器}
\ollabel{fig:adder}
\end{figure}
\end{ex}

\begin{prob}
追踪 \olref[tur][mac][una]{ex:adder} 中机器在输入为~$\tuple{3,2}$ 时的运行过程。
如果该机器计算 $0+0$，会发生什么？
\end{prob}

\begin{explain}
在 \olref[rep]{ex:doubler} 中，我们曾给出一个图灵机的例子，
它以一串 $\TMstroke$ 作为输入，停机时纸带上留下数量为其两倍的 $\TMstroke$ 串——这就是加倍器机器。
然而，由于输出在翻倍的 $\TMstroke$ 块左侧包含了 $\TMblank$，
它实际上并没有计算函数 $f(x) = 2x$，尽管你可能曾以为它计算了。
我们将描述两种修正方法。
\end{explain}

\begin{ex}
\olref{fig:doubler-disc} 中的机器计算函数 $f(x) = 2x$。
\olref[rep]{ex:doubler} 中的机器会擦除输入，并为输入中的每个 $\TMstroke$ 在最右侧写下两个 $\TMstroke$；
而本机器则不同，它为输入中的每个 $\TMstroke$ 在右侧添加一个 $\TMstroke$。
它需要跟踪输入在哪里结束，因此它在输入和添加的 $\TMstroke$ 之间留了一个 $\TMblank$，
并在最后用 $\TMstroke$ 将其填补。
同时，为了“记住”当前处理到输入中的哪个位置，
我们临时将输入块中的一个 $\TMstroke$ 替换为 $\TMblank$。
  \begin{figure}
\[
\begin{tikzpicture}[->,>=stealth',shorten >=1pt,auto,node distance=2.8cm,
                    semithick]
  \tikzstyle{every state}=[fill=none,draw=black,text=black]

  \node[initial,state] (0)              {$q_0$};
  \node[state]         (1) [above of=0] {$q_1$};
  \node[state]         (2) [above of=1] {$q_2$};
  \node[state]         (3) [right of=2] {$q_3$};
  \node[state]         (4) [below of=3] {$q_4$};
  \node[state]         (5) [below of=4] {$q_5$};
  \node[state]         (6) [above right of=4] {$q_6$};
  \node[state]         (7) [below of=6] {$q_7$};
  \node[state]         (8) [below of=7] {$q_8$};

  \path (0) edge node {\TMtrans{\TMstroke}{\TMblank}{\TMright}} (1)
%    (0) edge node {\TMtrans{\TMblank}{\TMblank}{\TMstay}} (h)
    (1) edge [loop left] node {\TMtrans{\TMstroke}{\TMstroke}{\TMright}} (1)
      edge node {\TMtrans{\TMblank}{\TMblank}{\TMright}} (2)
    (2) edge [loop above] node {\TMtrans{\TMstroke}{\TMstroke}{\TMright}} (2)
        edge node {\TMtrans{\TMblank}{\TMstroke}{\TMleft}} (3)
    (3) edge [loop above] node {\TMtrans{\TMstroke}{\TMstroke}{\TMleft}} (3)
        edge node[left] {\TMtrans{\TMblank}{\TMblank}{\TMleft}} (4)
    (4) edge        node[left] {\TMtrans{\TMstroke}{\TMstroke}{\TMleft}} (5)
    (5) edge [loop below] node {\TMtrans{\TMstroke}{\TMstroke}{\TMleft}} (5)
        edge              node {\TMtrans{\TMblank}{\TMstroke}{\TMright}} (0)
    (4) edge      node[sloped] {\TMtrans{\TMblank}{\TMstroke}{\TMright}} (6)
    (6) edge              node {\TMtrans{\TMblank}{\TMstroke}{\TMright}} (7)
    (7) edge [loop right] node {\TMtrans{\TMstroke}{\TMstroke}{\TMright}} (7)
        edge              node {\TMtrans{\TMblank}{\TMblank}{\TMleft}} (8)
    (8) edge [loop right] node {\TMtrans{\TMstroke}{\TMblank}{\TMstay}} (8);
    \end{tikzpicture}
\]
\caption{计算 $f(x) = 2x$ 的机器}
\ollabel{fig:doubler-disc}
\end{figure}
\end{ex}

\begin{ex}\ollabel{ex:mover}
计算 $f(x) = 2x$ 的第二种可能方案是保留原来的加倍器机器，
但在末尾添加状态和指令，将翻倍的 $\TMstroke$ 块移动到纸带的最左端。
\olref{fig:mover} 中的机器正是完成这最后一部分工作：
它从一段 $\TMblank$ 后接一段 $\TMstroke$（并且读写头位于 $\TMblank$ 块中的任意位置）的纸带开始，
逐个擦除 $\TMstroke$ 并将它们写到纸带的开头。
为了能够判断何时处理完毕，它首先用 $\TMendtape$ 符号标记 $\TMstroke$ 块的末尾，并在最后删除该符号。
我们将状态编号从 $q_6$ 开始，以便将其添加到加倍器机器上。
你只需要额外添加一条指令 $\delta(q_0,
\TMblank) = \tuple{q_6,\TMblank,\TMstay}$，即一条从 $q_0$ 到 $q_6$、标为 $\TMtrans{\TMblank}{\TMblank}{\TMstay}$ 的箭头。
（有一个微妙的问题：组合后的机器对输入 $x=0$ 不起作用。我们将此留作练习。）
\begin{figure}
  \[
  \begin{tikzpicture}[->,>=stealth',shorten >=1pt,auto,node distance=2.8cm,
                      semithick]
    \tikzstyle{every state}=[fill=none,draw=black,text=black]
    \node[initial,state] (6)              {$q_6$};
    \node[state]         (7) [right of=6] {$q_7$};
    \node[state]         (8) [right of=7] {$q_8$};
    \node[state]         (9) [below of=8] {$q_9$};
    \node[state]         (10) [left of=9] {$q_{10}$};
    \node[state]         (11) [left of=10]  {$q_{11}$};
    \node[state]         (12) [below of=11]  {$q_{12}$};
    \node[state]         (13) [right of=12] {$q_{13}$};
    \node[state]         (14) [right of=13] {$q_{14}$};
    \path
    (6)  edge node {\TMtrans{\TMblank}{\TMblank}{\TMright}} (7)
    (7)  edge [loop above] node {\TMtrans{\TMblank}{\TMblank}{\TMright}} (7)
         edge node {\TMtrans{\TMstroke}{\TMstroke}{\TMright}} (8)
    (8)  edge [loop above] node {\TMtrans{\TMstroke}{\TMstroke}{\TMright}} (8)
         edge node {\TMtrans{\TMblank}{\TMendtape}{\TMleft}} (9)
    (9)  edge [loop right] node {\TMtrans{\TMstroke}{\TMstroke}{\TMleft}} (9)
         edge node {\TMtrans{\TMblank}{\TMblank}{\TMright}} (10)
    (10) edge node[above] {\TMtrans{\TMstroke}{\TMblank}{\TMleft}} (11)
    (11) edge [loop above] node {\TMtrans{\TMblank}{\TMblank}{\TMleft}} (11)
         edge node[left] {\begin{tabular}{@{}l@{}}
          \TMtrans{\TMendtape}{\TMendtape}{\TMright}\\
          \TMtrans{\TMstroke}{\TMstroke}{\TMright}
         \end{tabular}} (12)
    (12) edge node[] {\TMtrans{\TMblank}{\TMstroke}{\TMright}} (13)
    (13) edge [loop below] node {\TMtrans{\TMblank}{\TMblank}{\TMright}} (13)
         edge node[sloped] {\TMtrans{\TMstroke}{\TMblank}{\TMleft}} (11)
    (13) edge node {\TMtrans{\TMendtape}{\TMblank}{\TMstay}} (14);
  \end{tikzpicture}
  \]
  \caption{将 $\TMstroke$ 块移到左侧}
  \ollabel{fig:mover}
\end{figure}
\end{ex}

\begin{prob}
在 \olref[tur][mac][una]{ex:mover} 中，我们描述了一台由
\olref[tur][mac][una]{fig:doubler-disc} 中的加倍器机器和
\olref[tur][mac][una]{fig:mover} 中的移动器机器组合而成的机器。
如果你以输入 $x=0$（即空白纸带）启动这台组合机器，会发生什么？
你如何修复此机器，使得在这种情况下机器停机时的输出为 $2x=0$？
（你应该能通过添加一个状态和一个转移来实现。）
\end{prob}

\begin{prob}
\emph{减法：} 设计一个图灵机，使得当输入为两个非空的 $\TMstroke$ 串，长度分别为 $n$ 和 $m$ 且 $n > m$ 时，它能计算函数 $f(n,m) = n - m$。
\end{prob}

\begin{prob}
\emph{相等：} 设计一个图灵机来计算以下函数：
\[
\fn{equality}(n,m) = 
\begin{cases}
  \text{1} & \text{如果~$n = m$} \\
  \text{0} & \text{如果~$n \neq m$}
\end{cases}
\]
其中 $n$ 和 $m \in \PosInt$。
\end{prob}

\begin{prob}
设计一个图灵机来计算函数 $\min(x,y)$，其中 $x$ 和 $y$ 是正整数，在纸带上由两串被一个 $\TMblank$ 分隔的 $\TMstroke$ 表示。你可以在机器的字母表中使用额外的符号。

函数 $\min$ 从其参数中选取最小的值，因此 $\min(3,5)=3$，$\min(20,16)=16$，$\min(4,4)=4$，依此类推。
\end{prob}

\begin{defn}
  图灵机~$M$ 计算偏函数 $f\colon \Nat^k \pto \Nat$，当且仅当满足：
  \begin{enumerate}
    \item 如果 $f(n_1, \dots, n_k) = m$，则 $M$ 在输入 $\TMstroke^{n_1}\concat\TMblank\concat
    \dots \concat\TMblank\concat\TMstroke^{n_k}$ 上停机，且输出为 $\TMstroke^{m}$。
    \item 如果 $f(n_1, \dots, n_k)$ 无定义，则 $M$ 要么根本不停机，要么停机时的输出不是单独的一块 $\TMstroke$。
  \end{enumerate}
\end{defn}

\end{document}